\section{相关工作}\label{sec:related}
协议逻辑已经说明，读取共享状态能够获得关于协议进展的知识。GPS 使用单地址协议、幽灵状态与资源转移推理弱内存程序，并以循环缓冲展示资源复用~\cite{gps2014}；iGPS 的读规则允许从当前已知协议状态前进到某个后继状态，单写者协议进一步支持常见更新纪律~\cite{igps2017}。因此“同步能传递不可逆知识”是本文的推理基础。本文进一步研究这种共同知识对\emph{所有尚未选择的候选}产生什么图结构限制：在候选过滤饱和后，稳定历史只保留每地址的少量极大写，再连接完整原始支持。缓冲窗口来自双向反馈；屏障的来源下界还利用了到达位的真实重置，不依赖完整交替日志的末端计数。

内存关系分析为候选与必然事实提供了成熟抽象。Haas 等人的 may/must 过滤框架允许程序信息、关系方程与模型公理双向传播~\cite{haas2023}。其具体过滤语义定理与可靠抽象过滤需要区分：可靠简化保留合法执行，但近似域本身并不等于精确支持。本文的 $R_\forall$ 可以作为这类分析中的具体规则；新结构结论需要额外证明 R1/R2 饱和、稳定极大写及 RF／未决 CO／覆盖边的共同边界，而不能从 may/must 可靠性直接读出 $K$ 的参数界。CAAT 将一致性作为可用于 SMT 的理论，并研究关系固定点计算和冲突解释~\cite{haas2022}。其 §4.1 讨论删除边后复用固定点的困难。本文通过预编译固定输入，对路径标签条件化来支持撤销；这里撤销的是原始假设，事件结构保持固定。

给定 RF 后补全修改序，与同时保留所有 RF／CO 选择是相关但不同的任务。Tunç 等人的 reads-from consistency checking 接收事件、程序顺序和已经选择的 RF，检查是否存在合法修改序，并对多种 C11 风格模型给出高效算法~\cite{tunc2023}。这项工作已经处理未知 CO。本文需要将不同 RF 选择共同保留，并回答任意附加 CO 条件后的全部支持，因此最终表示不能把 CO 仅作为无须查询的存在见证消去。

Conrado 与 Pavlogiannis 的 MSO 框架研究已选 RF 图的结构和弱内存可定义性~\cite{conrado2026}。其 Lemma~1 约束 SC 的 RF 图树宽，Lemma~4 把高树宽联系到任意 $\SB\cup\rf$ 拓扑序中的大写距离；文中也明确展示长 RF 边可以与低树宽并存。其 Theorem~9 和 Lemma~10 给出 RA 等模型在 RF 图上的 MSO 表达，结合结构结果得到验证与 RF 鲁棒性算法。本文控制的对象还包括全部 RF 候选、未决 CO 与固定 CO 覆盖边，所用布局先于所有完成确定。第~\ref{sec:frontier}节的 $K_{h,h}$ 例子说明为何逐完成结构界不能直接推出这一结论。稳定写代表进一步把共同因果关系转为可构造的候选接口界；长边存在但端点少这一观察本身不是贡献边界。

复杂度结果解释了为何还需要协议条件。Chakraborty 等人证明，RA 及其相关变体的一致性测试即使在线程数、地址数和值域有界且没有 RMW 时，仍为 NP 完全~\cite{chakraborty2024}。Govind 等人进一步研究每地址写者数：单写者的 RA／SRA／WRA 一致性测试可多项式求解，而允许两个写者时已有 NP 完全性~\cite{govind2026}。这些结论针对各自的观测测试问题，不直接成为本文协议子类的下界。特别是单写者缓冲的实验价值在于重复条件支持及实际跨阶段关联，不能表述成把一个已知困难的单次判定问题变为易解。多写者屏障则说明前沿参数可以在写者数大于一时仍然受控。

关系自动机与知识编译提供后端基础。Kater 用关系代数、自动机及闭合游走检查处理内存模型元理论和一致性~\cite{kater2023}；本文沿用正则关系监视的思路。路径表示上的条件化、投影与查询能力属于知识编译的标准研究对象~\cite{darwiche2002}，由平滑 DNNF 构造支持传播编码也已有系统结果~\cite{kucera2021}。本文特别保持原始变量投影，使用同一历史的普通／一标记摘要建立可接受后缀等价，再保留路径标签。通用编译步骤不是单独的创新主张。d4 等编译器已经利用动态分解和组件缓存~\cite{lagniez2017}；本文静态 BD 对照用于识别局部分解对收益的解释范围，没有与这些成熟编译器在相同任务上作性能排名。
